{\bfseries Work in progress!!}

\subsection*{Introduction}

Ge\+N-\/\+Foam has been developed for the steady-\/state and transient analysis of reactors featuring pin-\/type, plate-\/type, or liquid fuel (viz., Molten Salt Reactors). It includes sub-\/solvers for neutronics, single-\/ and two-\/phase thermal-\/hydraulics, and thermal-\/mechanics, with the choice of the physics to solve that can be made at runtime. Three different meshes are employed for neutronics, thermal-\/hydraulics, and thermal-\/mechanics. The selection of the physics to solve is done in {\itshape system/control\+Dict}

 {\bfseries The {\itshape control\+Dict} dictionary}

The {\itshape control\+Dict} is an extended version of the one that is normally used in other Open\+F\+O\+AM solvers. Compared to a standard Open\+F\+O\+AM control\+Dict, it includes several keywords that allow to select\+: 
\begin{DoxyItemize}
\item Which physics to solve via the keywords {\itshape solve\+Fluid\+Mechanics}, {\itshape solve\+Energy}, {\itshape solve\+Neutronics}, solve\+Thermal\+Mechanics 
\item The type of reactor via the keyword {\itshape liquid\+Fuel} 
\item The options for time stepping via the keywords {\itshape adjust\+Time\+Step}, {\itshape max\+DeltaT}, {\itshape max\+Co} and {\itshape max\+Power\+Variation}, as well as {\itshape max\+Co\+Two\+Phase} and {\itshape margin\+To\+Phase\+Change} for two-\/phase flow simulations. 
\item An option for mesh manipulation called {\itshape remove\+Baffles}. This flag is not mandatory and allows to create a ghost thermal-\/hydraulics mesh without baffles.\+This ghost mesh allows for better mesh-\/to-\/mesh projections between different physics. W\+A\+R\+N\+I\+NG\+: parallel execution not tested. 
\end{DoxyItemize}Fairly complete examples of {\itshape control\+Dict} for single-\/phase flow can be found in \href{https://gitlab.com/foam-for-nuclear/GeN-Foam/-/blob/master/Tutorials/2D_FFTF/rootCase/system/controlDict}{\tt 2\+D\+\_\+\+F\+F\+TF} and \href{https://gitlab.com/foam-for-nuclear/GeN-Foam/-/blob/master/Tutorials/3D_SmallESFR/rootCase/system/controlDict}{\tt 3\+D\+\_\+\+Small\+E\+S\+FR}, while an explanation of the two-\/phase flow options can be found in \href{https://gitlab.com/foam-for-nuclear/GeN-Foam/-/blob/master/Tutorials/1D_boiling/system/controlDict}{\tt 1\+D\+\_\+boiling}. 

\subsection*{Coupling logic}

The coupling between physics is achieved by projecting coupling variables from the mesh they are calculated, to the mesh they need to be used. The following figure shows the overall logic behind the coupling.



The thermal-\/hydraulics sub-\/solver operates on a certain computational domain $\Omega_{TH}$. Given an input volumetric power density $q$, the thermal-\/hydraulics sub-\/solver is tasked with predicting the resulting fluid temperature $T$, density $\rho$ and velocity $u$ fields, as well as relevant structure temperature fields $T_s$. For a two-\/phase treatment, the fields $\rho$, $T$, $u$ consist of mass-\/weighed mixture values. The velocity field $u$ is used for coupling only when simulating M\+S\+Rs to advect the precursors. The field $q$ is the volumetric fuel power density and it can pertain either to a sub-\/scale structure (typically, the fuel rods) or the fluid itself (i.\+e. the liquid fuel in M\+S\+Rs), depending on the system under investigation. The symbol $T_s$ collectively denotes the temperature fields of the structures, which can range from the fuel and cladding of a nuclear fuel pin to control rod drivelines, wrappers, the diagrid, etc. This entirely depends on what the structure thermal models are supposed to represent in the cell zones where they have been defined.

The neutronics sub-\/solver operates on a computational domain $\Omega_N$ and is tasked with predicting the volumetric fuel power density $q$ for varying coupling fields. Not all of these fields are always used, depending on the selected type of neutronics treatment. In general terms, the diffusion, $S_N$, $SP_3$ treatments are capable of modeling reactivity feedbacks from\+: coolant temperature $T$ and density $\rho$, average fuel and cladding temperatures collectively denoted with $T_s$, fuel axial displacement and core radial displacement collectively denoted as $d$. As long as a parametrization of the macroscopic cross-\/sections against these quantities is provided, these feedbacks can be resolved. The feedback reactivities of the point-\/kinetics solver are described by standard feedback coefficients.

The thermal-\/mechanics sub-\/solver operates on a computational domain $\Omega_{TM}$ and is tasked with predicting an overall displacement field that can be used to deform all the computational domains $\Omega_{TH}$, $\Omega_{N}$, $\Omega_{TM}$. The displacement field is decomposed into fuel axial displacement and core radial displacement fields collectively denoted as $d$, which are passed to the neutronics to model expansion-\/related feedbacks.

\subsection*{Prioritization in coupling variables}

This is one of the most complex aspects of multi-\/physics solvers like Ge\+N-\/\+Foam. First of all, one should understand that, in multi-\/physics simulations, the coupling variables $q$, $T$, $\rho$, $u$, $T_s$ and $d$ must exist both in their original form in the region (mesh) where they are created, and in their projected form in the region (mesh) where they are used. With that in mind, the question immediately arises about the origin of variables when mutiple sub-\/solver are activated.

In Ge\+N-\/\+Foam, for nearly all the coupling fields, a sub-\/solver will look at the original form of the coupling variable. For instance, the neutronic sub-\/solver will read the temperature fiel of the fuel from the thermal-\/hydraulics sub-\/solver (i.\+e., from the {\itshape fluid\+Region}). This is physically intutitive ans allows Ge\+N-\/\+Foam to behave as expected in sequential runs. For instance, one will be able to run a thermal-\/hydraulic calculations first, and then start from there and automatically use in a neutronic calculation the resulting temperature and density fields in order to properly parametrize the cross-\/sections.

However, a drawback exists for a user that wishes to execute a single-\/physics simulation, but to provide its own specific value to the coulpling field. The standard behaviour of Open\+F\+O\+AM will force this user to provide this value in the region where the field is normally calculated. Following the example above, a user wishing to run a neutronicscalculation with its own temperature and densities will have to generate a mesh for the {\itshape fluif\+Region}, and provide the desired temperatures and density fields in the 0 (or other {\itshape start\+Time}) folder of that region.

While the mentioned drawback is expected to negatively impact a minority of Ge\+N-\/\+Foam users, this is not the case for $q$ (i.\+e., the {\itshape power\+Density} field), for which it was decided to adopt a separate treatment. In this case, Ge\+N-\/\+Foam will take the {\itshape power\+Density} field from the {\itshape flui\+Region} when neutronics is switched off, while the neutronic sub-\/solver will project its own {\itshape power\+Density} field to the {\itshape neutro\+Region} when neutronics is switched on. This allows for stand-\/alone thermal-\/hydraulic simulations with a power density calculated e.\+g. using Serpent, without the need to create a mesh for the neutronics. At the same time, whrn neutronics is switched on, it will overwrite the {\itshape power\+Density} of the {\itshape fluid\+Region} with its own calculated value. More information about how to set the initial power\+Density in Ge\+N-\/\+Foam calculations is provide in the \char`\"{}\+Setting the $\ast$power\+Density$\ast$\char`\"{} box in \hyperlink{TH}{Thermal-\/hydraulics}.

Of course, a consistent and unified apparoach could be obtained for all couplig variables by introducing additional flags to let the user decide where the fields must be taken from, and projected to, but for the moment it was decided that this would add unnecessary complexity to code and input deck.

\subsection*{Coupling loop}

The coupling between physics is obtained via fixed-\/point iterations and the time derivatives are solved based on a first-\/order implicit Euler scheme. The figure below reports the overall coupling scheme. Within each time step, the outer iteration loop is used to couple the single-\/physics solution steps it encompasses. It can be performed a user-\/selected number of times or controlled by the convergence of the residuals of the slowest-\/converging physics.



The parameters for the coupling are set in {\itshape system/fv\+Solution}.

 {\bfseries The general {\itshape fv\+Solution} dictionary}

The general {\itshape fv\+Solution} dictionary is found under $\ast$/system/$\ast$ and allows to specify parameters related to the coupling among physics, and in particular\+: the type of coupling (implicit or explicit); and the parameters that affect the tightness pof the implicit coupling.

A commented {\itshape fv\+Solution} can be found in \href{https://gitlab.com/foam-for-nuclear/GeN-Foam/-/blob/master/Tutorials/3D_SmallESFR/rootCase/system/fvSolution}{\tt 3\+D\+\_\+\+Small\+E\+S\+FR}. 

© All rights reserved. E\+C\+O\+LE P\+O\+L\+Y\+T\+E\+C\+H\+N\+I\+Q\+UE F\+E\+D\+E\+R\+A\+LE DE L\+A\+U\+S\+A\+N\+NE, Switzerland, 2021 